\scanpage{40}
$U_\rho$ opère (sans \unsure{mérite|unité}!) sur $D_0$ et $D_0^{**}$, et que $D_0^{**}$ est un \uline{torseur} sous $U_\rho$\unsure{,} (trivialisé par $s_0\ldots$)

\begin{center}\rule{2.1cm}{0.35pt}\end{center}

Supposons maintenant les $\mu_i=2$ de sorte qu'on a une forme quadratique \uline{inv.} canonique $f$ sur $\unsure{E|\mathcal E}$, unique par la condition

\[
f(a_{-1})=0,\qquad f(e_0)=1,
\]

\deleted{\unsure{Je di|Jx de}} d'où $Q\subset X=\mathbb P^{n}(E)$.

On sait que la ps \unsure{réfle|réfl.} $\mathrm{id}+\underbrace{\lambda\pi'\otimes\pi}_{\pi'\otimes\lambda\pi}$ est \deleted{\unsure{une}} \unsure{orthogonal|orthogonale} sss \deleted{\unsure{si}}l'opération corr. dans $X$ conserve $Q$, et ceci signifie que

\[
\begin{aligned}
&\deleted{\unsure{\lambda}}\unsure{f(\deleted{\unsure{\psi f|\psi|\varphi}}\lambda\pi)|f(\pi)}\pi'+\psi(\lambda\pi)=0\quad\text{i.e.}\\
&\lambda\bigl[\lambda\underbrace{f(\pi)}_{-\widetilde\sigma\widetilde\rho}\pi'+\underbrace{\psi(\pi)}_{-\sigma\pi'}\bigr]=0\qquad\text{i.e.}\\[-2pt]
&\hspace{1.5em}\deleted{\unsure{\lambda f\pi'+{}}}
\end{aligned}
\]

\[
\lambda(\sigma+\lambda\widetilde\sigma\widetilde\rho)=0
\tag*{$(**)$}
\]
% archive sheet number: 39
